\documentclass{article}
\usepackage{graphicx}
\usepackage{subcaption}% Required for inserting images

\title{Cloud physics mandatory project}
\author{beateoevland }
\date{April 2026}

\begin{document}



\maketitle


\section{Problem 1}

\subsection*{a) Aerosol Distribution Modes from Ship Emissions}

Ship emissions are the result of combustion of fossile fuels, which contribute to the Fine Mode ($d < 2.5 \micro m$) of the aerosol distribution. There are two ways this happens: 

\subsubsection*{Primary Aerosols}

Ships emit unburned carbon directly into the atmosphere, known as black carbon (soot). These primary particles are extremely small, and fall into the fine mode. 

\subsubsection*{Secondary Aerosols}

Ship exhaust contains gases like $SO_2$, $NO_x$, and volatile organic compounds (VOC). In the atmosphere, sunlight triggers chemical reactions that change these gases. $SO_2$ can oxidize to $H_2SO_4$, which can then react with water and ammonia to form tiny particles called sulfate aerosols. These particles are very small and fall within the fine mode size range. 

\subsection*{b) How ship emissions modify clouds}

Ship exhaust can drastically alter the local cloud properties through processes known as Aerosol-Cloud Interactions (Indirect effects). 

\subsubsection*{Acting as a Cloud Condensation Nuclei (CCN)}

Secondary aerosols formed from ship emissions (sulfate salts in particular) are highly water-soluble. According to Köhler theory, dissolving a solute into a droplet lowers the equilibrium vapor pressure required for the droplet to grow. Therefore, these ship emissions serve as effective CCNs. 

\subsubsection*{Increased Albedo}

Ship emissions introduce a massive amount of CCNs into the marine boundary layer. As a result, the available condensed water is divided among many more particles, creating a cloud with a much higher concentration of smaller cloud droplets. Because smaller droplets have a larger total surface area, they scatter solar radiation more efficiently, significantly increasing the cloud's brightness (albedo).  

\subsubsection*{Increased lifetime}

Since the cloud droplets are smaller, it takes them much longer to collide, coalesce, and grow heavy enough to fall. This suppresses drizzle and precipitation, which increases the cloud's height, thickness, and lifetime. 

\subsection*{c) Cloud Formation in Aircraft Contrails}

Formation of contrails behind an aircraft is an example of isobaric adiabatic mixing. At high altitudes, engines burn fuel and release a plume of exhaust that is both very hot and enriched with water vapor. The surrounding air is dry and extremely cold. When the hot, moist exhaust mixes with the dry, cold surrounding air, the temperature and the partial pressure of water vapor of the mixture adjust as a linear combination of the two air masses. 

Because the relationship between temperature and saturation water vapor-pressure is non-linear (described by the Clausius-Clapeyron relation), the mixed air can temporarily become supersaturated with respect to water. When this happens, the excess water vapor rapidly condenses or freezes onto particles in the exhaust, forming ice crystals. These ice crystals are visible as cloud trails behind the aircraft, known as contrails. 

\section*{Problem 2}

\subsection*{b) Dry and moist adiabat}

The dry and moist adiabats differ because in saturated air, condensation releases latent heat. This reduces the cooling rate compared to unsaturated air. In the Stavanger sounding, the difference between dry and moist adiabatic ascent is small near the surface because the lifting condensation level is low. However, the difference increases with height, where the moist parcel becomes significantly warmer than the dry parcel.


\subsection*{c) Vapor pressure and saturation vapor pressure}

The saturation vapor pressure is calculated using the Clausius–Clapeyron equation:

\[
e_s(T) = e_0 \exp\left( \frac{L_v}{R} \left( \frac{1}{T_0} - \frac{1}{T} \right) \right)
\]

where:
\begin{itemize}
\item $L_v = 40660 \ \text{J mol}^{-1}$
\item $R = 8.31 \ \text{J mol}^{-1} \text{K}^{-1}$
\item $e_0 = 6.11 \ \text{hPa}$
\item $T_0 = 273.15 \ \text{K}$
\end{itemize}

For the surface temperature $T = 8.6^\circ\text{C} = 281.75 \ \text{K}$:

\[
e_s = 6.11 \exp\left( \frac{40660}{8.31} \left( \frac{1}{273.15} - \frac{1}{281.75} \right) \right)
\]

\[
e_s \approx 10.56 \ \text{hPa}
\]

The actual vapor pressure is calculated from the dew point temperature:

\[
e = e_s(T_d)
\]

For $T_d = 4.1^\circ\text{C} = 277.25 \ \text{K}$:

\[
e \approx 7.96 \ \text{hPa}
\]


\subsection*{d) Dew point and LCL}

The dew point temperature at the surface is:

\[
T_d = 4.1^\circ\text{C}
\]

The LCL temperature is calculated using Bolton's approximation:

\[
T_{LCL} = \left( \frac{1}{\frac{1}{T_d - 56} + \frac{\ln(T/T_d)}{800}} \right) + 56
\]

where temperatures are in Kelvin.

\[
T = 281.75 \ \text{K}, \quad T_d = 277.25 \ \text{K}
\]

\[
T_{LCL} \approx 276.3 \ \text{K} = 3.1^\circ\text{C}
\]

The potential temperature is:

\[
\theta = T \left( \frac{1000}{p} \right)^{0.2854}
\]

\[
\theta \approx 281.75 \left( \frac{1000}{1013} \right)^{0.2854}
\]

The LCL pressure is then:

\[
p_{LCL} = 1000 \left( \frac{T_{LCL}}{\theta} \right)^{\frac{1}{0.2854}}
\]

\[
p_{LCL} \approx 946 \ \text{hPa}
\]

From interpolation of the sounding:

\[
z_{LCL} \approx 600 \ \text{m}
\]

\subsection*{e) Virtual temperature}

The mixing ratio is calculated as:

\[
w = \frac{0.622 \, e}{p - e}
\]

\[
w = \frac{0.622 \times 7.96}{1013 - 7.96}
\]

\[
w \approx 0.00493 \ \text{kg kg}^{-1}
\]

The virtual temperature is:

\[
T_v = T (1 + 0.61 w)
\]

\[
T_v = 281.75 (1 + 0.61 \times 0.00493)
\]

\[
T_v \approx 282.6 \ \text{K} = 9.45^\circ\text{C}
\]

Thus:

\[
T_v - T \approx 0.85^\circ\text{C}
\]

The virtual temperature is higher than the actual temperature because moist air is less dense than dry air. This affects buoyancy, since buoyancy depends on density differences between the air parcel and the surrounding environment.

\subsection*{f) Vertical profiles}

The temperature decreases with height throughout most of the troposphere. Near the surface, the temperature and dew point are relatively close, indicating moist air. Around 925--855 hPa, the temperature and dew point nearly coincide, suggesting a saturated layer where clouds may form.

Above this level, the dew point decreases more rapidly than the temperature, indicating much drier air. This dry layer limits vertical cloud development. A lifted parcel follows a dry adiabat below the LCL and a moist adiabat above it. The profile does not indicate strong instability, suggesting limited convective activity.

\subsection*{g) Cloud formation}

Clouds are expected to form where temperature and dew point are nearly equal:

\[
900 \text{--} 850 \ \text{hPa}
\]

This corresponds to approximately:

\[
z \approx 800 \text{--} 1400 \ \text{m}
\]

The cloud type is likely stratiform (stratus or stratocumulus), since the atmosphere shows a stable structure with limited vertical development.

A dry layer above the cloud layer inhibits deep convection, so significant precipitation is unlikely. Some light drizzle may occur if the cloud layer becomes sufficiently thick.

\subsection*{h) Vertical velocity from CAPE}

The maximum vertical velocity can be estimated from CAPE:

\[
w_{\text{max}} = \sqrt{2 \cdot \text{CAPE}}
\]

For CAPE = 2000 \ \text{J kg}^{-1}:

\[
w_{\text{max}} = \sqrt{2 \times 2000}
\]

\[
w_{\text{max}} \approx 63.2 \ \text{m s}^{-1}
\]

This value represents an ideal upper limit. In reality, vertical velocities are lower due to entrainment, turbulence, drag, water loading, and phase changes of water.


\section{Problem 3}

\subsection*{a)}

\subsubsection*{2)}

The CM1 (Cloud Model 1) is a three-dimensional, nonhydrostatic numerical model designed to simulate atmospheric flows at cloud-resolving scales. It solves the compressible Navier–Stokes equations using finite-difference methods on a structured grid.

The model includes representations of key physical processes such as:
- Atmospheric dynamics (advection, pressure gradients, buoyancy)
- Microphysics (formation and evolution of hydrometeors)
- Turbulence and subgrid-scale mixing
- Surface fluxes and boundary layer processes
- Optional radiation and Coriolis effects

CM1 is highly flexible and is commonly used in idealized simulations to study deep convection, supercells, squall lines, and boundary layer processes. Its main purpose is to improve understanding of storm dynamics, cloud processes, and interactions between thermodynamics and flow fields.

\subsubsection*{4)}

The Kessler scheme is a simple warm-rain microphysics parameterization that includes only liquid-phase processes. It represents cloud water (qc) and rainwater (qr), with parameterizations for condensation, evaporation, and conversion of cloud water into rain.

The scheme does not include ice-phase processes such as freezing, deposition, or graupel formation. As a result, it is computationally efficient but limited in its ability to represent upper-level cloud structure and mixed-phase processes in deep convection.

The Morrison scheme is a more advanced two-moment microphysics parameterization that predicts both the mass mixing ratio and number concentration of hydrometeors. It includes multiple parameters such as cloud water, rain, ice crystals, snow, and graupel.

The scheme represents a wide range of microphysical processes, including condensation, evaporation, freezing, melting, deposition, aggregation, and riming. This allows for a more realistic simulation of mixed-phase and ice processes in convective storms, particularly in the upper troposphere.

Compared to the Kessler scheme, the Morrison scheme provides a more detailed and physically realistic representation of cloud and precipitation processes, but at a higher computational cost.

\subsection*{b)}

\subsubsection*{1)}

The horizontal grid spacing is 1 km in both x and y directions, while the vertical grid spacing is 500 m. The model time step is 6 seconds, and the total simulation time is 14400 seconds (4 hours).

\subsubsection*{2)}

The horizontal domain is 120 km × 120 km, with 120 grid points in each direction. The vertical extent of the model is 18 km and 40 grid points.

\subsubsection*{3)}

\subsubsection*{6)}
Cloud phases: 

\title{Initial phase (timestep 0-3):}

At the beginning of the formation of the cell we have a weak convective cell. A localized updraft is developing near the surface. The cloud water (qc) is shallow and compact, and hasn't yet developed vertically, meaning convection is still weak. In this stage, warm air has begun to rise, but it isn't yet a storm. The flow field (wind and vertical velocity) shows a weaker updraft. There is very little graupel and cloud ice formed at this stage. 
 \\


\title{Growth phase (4-7):}

In this stage, we see rapid vertical development of the cloud water, which was shallower before. Cloud ice and graupel are also rapidly developing, and we can see a tall convective cloud column forming. The ice cloud is developing an anvil shape and spreading aloft. In this stage we have increasing buoyancy and latent heat release from the phase changes in the rising air. \\

\title{Supercell stage (8-11):}

The cloud has a well-defined core. 
Here we have a well-defined anvil cloud spreading horizontally at upper levels. Here we also see multiple updrafts developing, and see significant amounts of cloud water, cloud ice, and graupel. We also see precipitation forming at the surface. With each timestep, the cloud gets more asymmetrical. At this stage, there is both buoyancy-driven updraft and downdraft driven by precipitation. 

\title{Later stage (12-16):}

The storm starts to die out. We get smaller cores and a multicellular system. The anvil becomes broader, but less active. We have reduced buoyancy due to cold pools at the surface. There is more precipitation at the surface. 

Vertical and horizontal slices. The vertical slices are taken around the core of the storm (x=48, y=48). The horizontal slices are taken at 500m, 2000m, 5000, and 10000m, although the dataset seems to start at 250m, so there might be an offset of 250m. 

Potential temperature profiles: 

The vertical profiles at t=06 show that $\theta$ increases with height, indicating stable stratification. We do see some weak, localized perturbations and a slight warming in the midsection, indicating early latent heat release. At t=10, there is a clear warm anomaly in the mid- and left sections, which passes through the storm. Compared to t=06, we see colder temperatures forming near the surface. At t=10, there is a clear warm core. The difference between t=06 and t=10 shows that the updraft has strengthened, more latent heat has been released, and a mature supercell structure is forming. 

The horizontal profiles at t=09 show mostly uniform $\theta$ at low height and some cooler anomalies near the storm center. At higher altitudes (especially 5000m), we can see more clearly warm anomalies, which corresponds to the updraft region. At 10000m, the temperature is much smoother and higher than before. There are fewer localized variations. 

All the temperature profiles show a potential temperature that increases with height, which is a stable environment. Around the storm center there is warm air higher up due to updraft and latent heat release, and cold formations near the surface due to evaporation and downdraft. \\

Wind profiles: 

From the vertical w profile, we see a strong red area indicating an updraft and a strong blue area indicating a downdraft. It's localized in a column, where our supercell core is. The horizontal profile for w shows a weaker, more spread-out signal, with a few downdraft signatures. Here we see the cold pools forming, as discussed in our temperature profiles. At higher levels, the vertical motion is stronger, with updrafts and downdrafts more noticeable. This is situated at the core of our storm. \\

The uinterp horizontal profiles show strong gradients and a change from positive to negative flow. This is the inflow and outflow of the storm. The vinterp horizontal profile shows wind shear where the air flows differently on each side of the updraft. 

\subsection*{c)}

\subsubsection*{Cloud evolution: }

In the Kessler scheme, we have only cloud water (no ice) and precipitation at the ground. \\

\title{Initial phase (timestep 0-3):}

Here, the cloud is still weak and localized. A column is forming vertically, and there is very little precipitation at the ground. \\

\title{Growth phase (4-7):}

Updraft strengthens, and multiple cores form. The cloud is beginning to spread horizontally, with slight signs of an anvil shape. There is a lot more precipitation at the surface. Structure is messy and not organized. \\ 

\title{Main stage (8-11):}

Cloud has spread a lot. Multiple vertical cores, but very fragmented. Strong precipitation areas (red areas). We don't see any clear rotation or dominant updraft. The system is multicellular, but not really supercellular. Precipitation is spreading, which is weakening the updrafts. \\

\title{Later stage (12-16):}

The cloud is diffusing, widening more, less coherent. Surface precipitation remains strong and widespread. The system is mostly dominated by precipitation and downdrafts. No sustained updrafts. \\

\title{Summary}

We don't see a clear supercell structure; there is no single rotating updraft and no clear separation between updraft and precipitation. It's a disorganized multicell system. Because of no ice and no graupel, we see no latent heat release from freezing and no strong upper-level processes. This results in weaker updrafts and a less organized, shorter-lived convective system. 

\subsubsection*{Horizontal slices, t=09}

\title{Potential temperature:}

At 500m, there are mostly uniform values, but slight localized variations near the storm region. So the boundary layer remains relatively well mixed, and only weak cooling occurs near the downdraft and precipitation areas. At 2000m, more structure appears with slight gradients near the storm region. This indicates air is being lifted  (we have cooler gradients). \\

\title{Wind, u and v:}

The general pattern at both heights, for both wind components, shows clear spatial variability, with regions of stronger inflow and weaker outflow. At 500m the winds are smoother, and you can see them feeding the storm region, towards the updraft region. At 2000m, the winds become more irregular and spread out. There are both updrafts and downdrafts, and the winds are being disrupted and deflected. 

\title{Vertical wind, w, at t=09 and t=06:}



\end{document}
